\chapter{Background}
\label{sec:background}

This chapter introduces the concepts the rest of the document relies
on: invasive alien plants and \emph{Acacia dealbata} in Portugal;
satellite imagery and how a species is mapped from it; Earth
Observation foundation models, in particular TESSERA; citizen-science
observations and their properties; and the two methodological
problems that follow from combining them, learning from presences
only and validating a model over space.

\section{Invasive Alien Plants}
\label{sec:background-iap}

Invasive alien plants are species introduced outside their native
range that establish and spread without further human
intervention.\todo{Source for the definition; none in the reading list
yet.} Portugal regulates them through Decree-Law No.~92/2019, which
lists the species considered invasive in the national territory;
several \emph{Acacia} species are on that list, and the national
forest inventory follows the same list when it records invasive
species in its field plots \cite{icnf2019}.

\subsection{\emph{Acacia dealbata} in Portugal}
\label{sec:background-iap-acacia}

\emph{Acacia dealbata} Link, the silver wattle, is a tree native to
south-eastern Australia.\todo{Source for the introduction history and
current extent in Portugal; none in the reading list.} It forms dense
stands that exclude native vegetation and also occurs scattered
inside stands of other species, where the forest inventory finds it
far more often than as stands of its own, as the figures given in
the motivation (Section~\ref{sec:introduction-motivation}) show.

Two traits of the species matter for this work. The first is its
flowering period. \emph{A.~dealbata} flowers in late winter, when most
of the surrounding vegetation does not, and the studies that have
mapped it from satellite imagery in Iberia found their best signal in
that period: with March imagery in central-eastern Portugal
\cite{martins2016} and with the spectral differences between the
pre-flowering, flowering and post-flowering months in Galicia
\cite{domingo2023}. The second trait is its response to fire. The
species keeps a persistent seed bank: its seeds have a physical
dormancy that heat breaks, and scarified seeds germinate at almost all
temperatures \cite{dessi2021}. In eight stands of the Coimbra district
treated between 2018 and 2021, burning promoted germination, slashing
stimulated resprouting, and no treatment prevented regeneration within
three years \cite{riveiro2025}. A burnt stand is therefore not a
cleared one.

\subsection{Why mapping matters}
\label{sec:background-iap-mapping}

Control of an invasive plant starts from knowing where it is, and the
cost of control grows with the area it has already occupied. The
information available in Portugal at species level is coarse in time
and space. The forest inventory, described in the motivation
(Section~\ref{sec:introduction-motivation}), is a decennial sample
that locates no stand. Land-cover maps
such as the Carta de Uso e Ocupação do Solo (COS) have an
\emph{Acacia} forest class, but with a minimum mapping unit of one
hectare they do not resolve scattered presences, and they mix land use
with land cover, which makes them a weak reference for 10~m
classification \cite{fonte2020}. Recurrent large fires, concentrated
in the centre of the country, as the fire statistics in the
motivation (Section~\ref{sec:introduction-motivation}) show, keep
opening new ground
for a species that regenerates after burning. A method that produces a
species map cheaply, and can repeat it every year, is what the
MAP.INVADER project, of which this dissertation is part, aims at.

\section{Earth Observation}
\label{sec:background-eo}

Earth Observation is the measurement of the Earth's surface from
satellites or aircraft. For vegetation mapping the relevant
instruments are passive optical sensors, which record the sunlight
reflected by the surface in several wavelength intervals, called
spectral bands. Vegetation reflects strongly in the near-infrared and
weakly in the red, so the relation between bands separates plants from
soil and water, and different plants, with different leaf structure
and phenology, differ in how that relation changes over the year.

\subsection{Sentinel-2 imagery and hand-designed features}
\label{sec:background-eo-sentinel2}

Sentinel-2 is the optical mission of the European Copernicus
programme. Two identical satellites carry a multispectral instrument
with thirteen bands from the visible to the short-wave infrared, at
10, 20 or 60~m depending on the band, and together revisit the same
location every five days at the equator and more often at higher
latitudes.\todo{Cite the ESA Sentinel-2 user guide or mission paper.}
The data are free and open. The Level-2A product gives surface
reflectance after atmospheric correction, together with a scene
classification layer that flags clouds, cloud shadows and snow, which
is the usual basis for masking unusable pixels \cite{ball2026}.

A single Sentinel-2 image is a grid of pixels, each holding one
reflectance value per band for one date. Because clouds hide the
surface on many dates, studies rarely use single images. They build
composites, for instance the median of all clear observations of a
month or a season, and derive indices from pairs of bands, such as the
normalised difference vegetation index. The \emph{Acacia} mapping
studies cited above are of this kind: monthly composites and index
differences between the flowering months \cite{domingo2023}, or one
image per month from December to May \cite{mouta2021}. The set of such
composites and indices for a pixel is its feature vector, and the
choice of which months and which indices to include is a modelling
decision made by hand.

\subsection{Pixel-based classification and reference data}
\label{sec:background-eo-classification}

Mapping a species from these data is a supervised classification
problem: a model learns, from pixels whose class is known, to assign a
class to every other pixel, and the map is the result of applying it
wall-to-wall. The pixels whose class is known are the reference data.
They serve two roles that must be kept apart: training the model, and
evaluating the map it produces on pixels the model has not seen.

Reference data can come from any source whose labels one is willing to trust. The conventional source is specialists, in two ways: field visits confirm the species at a location and give a point or a small plot, and visual interpretation of very-high-resolution imagery delimits whole stands as polygons, which is how the expert data used in this work were made, over 30~cm orthoimagery of 2023. Both are reliable, both are slow, and both are produced for one area and one year. Citizen-science observations are a second source, far cheaper and continuous but made of single points of uneven quality, as the section on citizen-science data
(Section~\ref{sec:background-cs}) describes; turning them into reference data usable at pixel level is one object of this dissertation. The scarcity of reliable reference data, not of imagery, is the limiting factor of supervised species mapping. In this work both sources are used as reference data, for training and for testing; what is asked is how far the cheap source can stand in for the expensive one.

\section{Earth Observation Foundation Models}
\label{sec:background-fm}

A foundation model is a large model pretrained once, on data without
labels, and reused for many tasks. This section describes what that
means for Earth Observation, where the reusable product is an
embedding per pixel, and then presents TESSERA, the model used in
this dissertation.

\subsection{Self-supervised pretraining and embeddings}
\label{sec:background-fm-embeddings}

The Sentinel-2 feature vector of
Section~\ref{sec:background-eo-sentinel2}, a hand-picked set of bands,
indices and seasonal composites, is designed anew for each study. Earth Observation foundation models
replace it with a representation learned from data: a large model
pretrained on a large volume of unlabelled imagery, by a
self-supervised objective that needs no reference data, and then
reused for many downstream tasks without being trained again. For
Earth Observation the reusable output is the \emph{embedding}: a
fixed-length vector per pixel, or per small patch, that summarises
what the model learned about that location from the full time series
of observations. Several such products are published as data rather
than as models, so that a user downloads the embeddings of an area and
a year and never runs the foundation model \cite{zhu2026}. AlphaEarth,
for example, provides 64-dimensional annual embeddings at 10~m for
2017 to 2024, computed on 1.28~km tiles with spatial self-attention,
so that a pixel's embedding also depends on its surroundings; only the
embeddings are released, not the model \cite{brown2025}. Downstream
use is then a light classifier on the embeddings. The two reasons to
prefer a frozen embedding to fine-tuning in this work are practical
and empirical: the models are either not released or restricted, and
the TESSERA paper itself found no benefit from fine-tuning its encoder
over using it frozen \cite{feng2025}.

\subsection{TESSERA}
\label{sec:background-fm-tessera}

TESSERA is the foundation model used in this dissertation. It is a
pixel-wise model: for each 10~m pixel it takes the year's Sentinel-2
multispectral time series and Sentinel-1 radar time series, built from
a random sample of valid dates, and produces a 128-dimensional
embedding for that year \cite{feng2025}. The embeddings cover 2017 to
2024 and are distributed through the GeoTessera library, so that the
embeddings of an area can be read without running the model. Two
properties follow from the pixel-wise design. First, the embedding of
a pixel does not depend on the pixels around it, unlike AlphaEarth;
the model's pretraining shuffles pixels globally to break spatial
correlation, and the landscape context is not encoded
\cite{feng2025,zhu2026}. Second, the embedding is annual: it
summarises the whole year's time series, and a disturbance inside the
year, such as a fire, changes it \cite{silva2026}. The model was
evaluated on tree species, crops, biomass and canopy height with light
classifiers on frozen embeddings, and in the comparison closest to
this work it beat a seasonal Sentinel-1 and Sentinel-2 baseline on 18
tree species in Trentino, reaching near-optimal weighted F1 with about
5\% of the training parcels \cite{ball2026}. No evaluation of TESSERA
on a species-occurrence or citizen-science task was found.

\section{Citizen Science Data}
\label{sec:background-cs}

Citizen science is the participation of non-specialists in the
collection of scientific data. In biodiversity it takes the form of
species observations: a person records that a species was seen at a
place and a date, usually with a photograph, on a platform that
aggregates the records. The volume and the continuity of these records
have no equivalent in professional surveys, which is why they are
being used as training and evaluation data for species maps
\cite{gillespie2024}.

\subsection{iNaturalist}
\label{sec:background-cs-inaturalist}

iNaturalist is the platform used in this work. Each observation holds
a date, coordinates, a photograph or sound, a taxon proposed by the
observer, and a positional accuracy, the radius in metres within which
the observer places the organism, which the observer may leave empty.
Other users can propose identifications. An observation reaches
\emph{research grade} when it has a date, a location and media, the
organism is wild, and the community agrees on the identification at
species level.\todo{Cite the
iNaturalist help page for the research-grade criteria.} Observations
whose coordinates the observer or the platform obscures have their
position randomised within a larger cell and cannot be used at pixel
level. Studies that use iNaturalist as training data keep only
research-grade observations and filter on positional accuracy, with
thresholds that depend on the resolution of the imagery: 10~m for
AlphaEarth's tree-genus evaluation \cite{brown2025} and 120~m for
256~m image patches \cite{gillespie2024}.

\subsection{Presence-only data, sampling bias and positional error}
\label{sec:background-cs-presence-only}

Three properties of these observations distinguish them from
reference data made by specialists. First, they are
\emph{presence-only}. A record says that the species was seen at a
point; the absence of a record says nothing, because no one may have
looked there. Confirmed absences, which a field survey produces, do
not exist in this kind of data \cite{phillips2009}. Second, the
sampling is opportunistic. Observers record where they go, so records
concentrate along roads, trails and accessible or disturbed places;
for invasive plants in Hawaii, iNaturalist records were biased
towards roads and disturbed vegetation while professional surveys
tended towards less accessible, native-dominated sites
\cite{dimson2023}. Third, the quality of each record is uneven. The
research-grade label is an agreement indicator: it measures how
reliable the identification is, not whether the position is accurate,
and agreement among contributors cannot replace a comparison with
reference data \cite{fonte2017}. The photograph's position may be
where the photographer stood rather than where the plant is,
``potentially some distance away'' \cite{fonte2017}, and at 10~m
resolution that distance is often a different pixel.

These properties do not make the data unusable; they define what has
to be done before using them. The first, the lack of confirmed
absences, shapes the learning problem, treated in the section on
learning and evaluating from presences only
(Section~\ref{sec:background-po-background}); the
second, the opportunistic sampling, shapes the validation, treated in
the section on spatial autocorrelation and block cross-validation
(Section~\ref{sec:background-po-spatial}); and the third, the uneven
quality of each record, shapes the preparation pipeline that turns
points into labelled pixels, described in the Proposed Work
(Chapter~\ref{sec:proposed-work}).

\section{Presence-Only Learning and Spatial Validation}
\label{sec:background-po}

Two of the properties of citizen-science data described above, the
lack of absences and the concentration of records where observers go,
raise two methodological questions that this section treats in turn:
how to train and evaluate a classifier from presences alone, and how
to validate it over space so that pixels next to the training data do
not inflate the result.

\subsection{Learning and evaluating from presences only}
\label{sec:background-po-background}

The missing class affects three things, taken here in order: the
choice of classifier, the sample of the landscape that stands in for
that class, and the evaluation of the result.

\subsubsection*{Two families of methods}

A standard classifier learns from examples of both classes.
Presence-only data give only one: the pixels where the species was
seen. There are two ways to train without the second class.

\emph{One-class methods} use the presences alone. They learn what a
presence pixel looks like and score every other pixel by how much it
resembles one. The one-class support vector machine does this by
fitting a boundary that encloses the presences in feature space as
tightly as possible, and scores a pixel by its distance to that
boundary. The isolation forest takes the opposite view: it builds
random trees that split the data, and a pixel that gets isolated in
few splits is unlike the presences and gets a low score.

\emph{Presence-background methods} create a second class by sampling
the landscape. A set of pixels is drawn at random over the whole study
area, called the \emph{background}, and a standard classifier is
trained to separate the presences from it. The background is not a set
of absences: a background pixel may well hold the species, since
nobody checked. The classifier's score therefore says how much a pixel
looks like the presences compared with the landscape in general, and
not how probable it is that the species is there \cite{phillips2006}.

This reading of the score holds for any classifier fitted with
background pixels: for MaxEnt, which was built on it, and equally for a
logistic regression or a random forest \cite{phillips2006}.
\citet{phillips2009} make the same point from the other side: the
background is an unlabelled sample, not a labelled one.

The distinction is often lost. A background pixel is one whose class
is unknown; a \emph{pseudo-absence} is a pixel the study decides to
treat as a true absence, and feeds to an ordinary
presence-versus-absence classifier whose output is then read as a
probability of presence. \citet{rapinel2021}, reviewing 136 one-class
vegetation studies, find that 42\% of them mix the two up, using one
and reading the result as if it were the other. The same review gives
two practical rules. The background should be drawn over the whole
study area, including the surroundings of the presences: if it is
taken only far from them, the classifier separates presences from very
different places and its accuracy is overestimated. And its size is
usually in the thousands, with 10\,000 points the most common choice.

Two other names appear in the literature, and both are
presence-background methods in disguise. \emph{MaxEnt} is the most
used model in species distribution modelling: it is fitted against
thousands of random background pixels and its output is read as a
relative index, so it differs from the methods above only in the form
of the model, not in what it learns from \cite{phillips2006}.
\emph{Positive-unlabelled learning} starts from the same
presence-versus-background classifier and rescales its score by the
fraction of presences that were recorded, to recover a probability;
when that fraction is the same everywhere, the rescaling does not
change the order of the pixels, and a rank-based metric sees the same
model \cite{elkan2008}. The real divide is therefore between one-class
methods, which use only the presences, and everything else, which also
uses the unlabelled landscape.

On that divide the evidence read for this work points one way. In
every remote-sensing comparison found, the methods that use the
landscape beat the one-class ones: a network trained on presences and
background against the one-class support vector machine, on aerial and
satellite scenes \cite{deng2018}; MaxEnt against the same machine
\cite{li2010}; a positive-unlabelled method against the biased and the
one-class support vector machines \cite{li2011}; and, in the one study
that puts the three families side by side, the biased support vector
machine first and the one-class one behind, with even the best method
performing poorly in many tasks \cite{mack2017}. \citet{elkan2008}
give the reason: one-class methods ``ignore useful information that is
present in the unlabeled examples''. The method-selection step of the
Proposed Work (Chapter~\ref{sec:proposed-work}) reaches the same
ordering on the data of this work.

\subsubsection*{Sampling bias in the background}

The random background carries a second problem. Because the presences
are biased towards accessible places and the background is not, the
model may learn where observers go rather than where the species
grows; \citet{phillips2009} show that a fitted model ``might be closer
to a model of survey effort than to a model of the true
distribution'', and propose drawing the background from the records
of other species collected in the same way, which share the bias. The
same paper shows that the gain of such a correction is visible only
against independent presence--absence data: scored against a random
background, the corrected model looks worse.

\subsubsection*{Evaluation without absences}

Evaluation inherits the same problem. Without absences, the metrics
are rank-based: how well the
model places the held-out presences above the landscape at large. One
form fixes the fraction of the study area the model is allowed to
predict as presence and measures the fraction of test presences that
fall inside it, so that two models are compared at equal predicted
area \cite{phillips2006,peterson2008}. The background then serves only
to measure that area: against a uniform random background, one minus
the specificity is almost identical to the proportion of the area
predicted present \cite{phillips2006}. An unrecorded presence at a
background point is never counted as an error, in training or in
evaluation. This is what makes ``the background is not absence'' true
in practice and not only in the reading of the model.

\subsection{Spatial autocorrelation and block cross-validation}
\label{sec:background-po-spatial}

Pixels close to each other are more alike than pixels far apart, in
their reflectance, in their embeddings and in their class. When
reference data are split at random into training and test sets, most
test pixels have a training pixel nearby, and the test measures how
well the model interpolates between neighbours rather than how well it
maps new ground. \citet{karasiak2022} measure the effect on Sentinel-2
forest classification: random cross-validation overestimated accuracy
by 8 to 14 points relative to distant sites, and the overestimate
disappeared when every test pixel was separated from the training set
by a distance set from the data, about 20~km in their case. The remedy
is to split by space rather than by pixel. In block cross-validation
the study area is divided into blocks, each block is held out in turn,
and all reference pixels inside it form the test set. Blocks alone
leave residual proximity across the block borders, so the papers that
compare designs recommend blocks combined with an exclusion buffer
around the test set, with the buffer set to the range of spatial
autocorrelation measured in the data, where underestimating the range
is the unsafe direction \cite{lerest2014,mahoney2023}. For
embedding-based models the design changes the conclusion: in a
comparison of TESSERA and AlphaEarth against Sentinel features for
old-growth forest, the advantage of the embeddings narrowed from 0.08
to 0.04 PR-AUC once a 10~km buffer separated training from test folds
\cite{ratsakatika2026}. The \emph{Acacia} mapping studies in Iberia
cited in the section on \emph{A.~dealbata} in Portugal
(Section~\ref{sec:background-iap-acacia}) all used random
splits \cite{martins2016,mouta2021,domingo2023}, so their accuracies
are not comparable with a spatially validated result.

\section{Overall Analysis}
\label{sec:background-overall}

\emph{A.~dealbata} is a species worth mapping and hard to keep mapped:
it spreads inside other stands, regenerates after fire, and the
official sources that describe it are decennial samples or coarse
land-cover maps. Sentinel-2 provides frequent, free, 10~m observations
from which the species can be told apart, mainly through its winter
flowering, but a supervised classifier needs reference pixels, and
those are the expensive part. Foundation models such as TESSERA move
the effort from designing features to collecting and aligning
reference data, by giving each pixel an embedding usable with a light
classifier and no fine-tuning. Citizen science offers reference data
at a scale and continuity that specialists cannot, but as presences
only, concentrated where observers go, with uneven positional quality.
Combining the two raises the two methodological questions of the
section on presence-only learning and spatial validation
(Section~\ref{sec:background-po}), which the Proposed Work
(Chapter~\ref{sec:proposed-work}) turns into the design: how to learn
and evaluate from presences without treating the rest of the landscape
as absence, and how to validate over space so that the result says
something about places the observers never visited.
